%%%%%%%%%%%%%%%%%%%%%%%%%%%%%%%%%%%%%%%%%%%%%%%%%%%%%%%%%%%%%%%%%%%%%%%%%%%
%
% Evaluation rubric for the Final Master's Thesis (TFM)
% EPS-UAH
%
% Structure and formatting based on Config/rubricaTFGTribunal.tex
% Contents reproduce Annex XI of the EPS-UAH TFM regulations.
% This file intentionally contains no person-dependent conditionals: the rubric
% wording and assessment criteria are fixed by the EPS-UAH regulation.
%
%%%%%%%%%%%%%%%%%%%%%%%%%%%%%%%%%%%%%%%%%%%%%%%%%%%%%%%%%%%%%%%%%%%%%%%%%%%

\begin{landscape}
\thispagestyle{plain}
\noindent \textbf{RÚBRICA PARA LA EVALUACIÓN DEL TFM}
\renewcommand{\arraystretch}{1.75}

\begin{table}[ht]
  \centering
  \scriptsize
  \begin{tabularx}{\linewidth}{|p{7cm}|p{9cm}|p{1cm}|p{1.5cm}|p{4.5cm}|}
    \cline{1-5}
    \textbf{Resultados de aprendizaje} & \textbf{Criterios de Evaluación} & \textbf{Peso} & \textbf{Calificación} & \textbf{Comentarios} \\ \cline{1-5}

    \multirow{6}{=}{%
      \begin{minipage}[t]{7cm}
        \textbf{Bloque 1: Proyecto}

        \vspace{1mm}
        RATFM2. Desarrollar proyectos relacionados con la ingeniería de telecomunicación con los estándares de calidad adecuados.

        \vspace{1mm}
        RATFM4. Definir todos los aspectos regulatorios de los proyectos en el campo específico de la ingeniería telecomunicación.

        \vspace{1mm}
        RATFM5. Integración de las competencias adquiridas en las enseñanzas en el desarrollo de un proyecto de ingeniería de
      \end{minipage}}
      & CE1. Integra las competencias adquiridas en las enseñanzas. & & & \multirow[t]{6}{=}{} \\ \cline{2-4}
      & CE2. Resuelve problemas con iniciativa, toma de decisiones y creatividad. & & & \\ \cline{2-4}
      & CE3. Tiene capacidad para el manejo de especificaciones, reglamentos y normas de obligado cumplimiento & & & \\ \cline{2-4}
      & CE4. Desarrolla proyectos relacionados con la ingeniería de telecomunicación con los estándares de calidad adecuados & & & \\ \cline{2-4}
      & CE5. Defiende un proyecto en el ámbito de la Ingeniería de telecomunicación. & & & \\ \cline{2-4}
      & \cellcolor[HTML]{C0C0C0}\textbf{Calificación máxima 5 puntos} & & & \\ \cline{1-5}

    \multirow{5}{=}{%
      \begin{minipage}[t]{7cm}
        \textbf{Bloque 2: Desarrollo del trabajo}

        \vspace{1mm}
        RATFM1.Interpretar adecuadamente las características de un proyecto de ingeniería de telecomunicación, comprenderlas y diseñar una aproximación al problema con creatividad e iniciativa propia.

        \vspace{1mm}
        RATFM6. Buscar y gestionar la información necesaria para dar respuestas a los retos planteados por un proyecto de ingeniería de telecomunicación.

        \vspace{1mm}
        RATFM7. Planificar las tareas a realizar para el desarrollo de un proyecto de ingeniería de telecomunicación.

        \vspace{1mm}
        RATFM10. Trabajar de forma autónoma, buscando
      \end{minipage}}
      & CE6. Busca la información necesaria para la realización del trabajo y la analiza para extraer aquella que le es de interés. & & & \multirow[t]{5}{=}{} \\ \cline{2-4}
      & CE7. Sintetiza la información obtenida de diferentes fuentes y los propios conocimientos en una visión global y estructurada del “estado del arte” del proyecto, referenciando correctamente la información obtenida de terceros. & & & \\ \cline{2-4}
      & CE8. Demuestra iniciativa, capacidad para tomar decisiones y capacidad de aprendizaje autónomo. & & & \\ \cline{2-4}
      & CE9. Planifica su trabajo, y es capaz de seguir la planificación y de analizar y justificar las posibles desviaciones sobre la planificación. & & & \\ \cline{2-4}
      & \cellcolor[HTML]{C0C0C0}\textbf{Calificación máxima 2 puntos} & & & \\ \cline{1-5}
  \end{tabularx}
\end{table}

\newpage
\thispagestyle{plain}

\begin{table}[ht]
  \centering
  \scriptsize
  \begin{tabularx}{\linewidth}{|p{7cm}|p{9cm}|p{1cm}|p{1.5cm}|p{4.5cm}|}
    \cline{1-5}
    \textbf{Resultados de aprendizaje} & \textbf{Criterios de Evaluación} & \textbf{Peso} & \textbf{Calificación} & \textbf{Comentarios} \\ \cline{1-5}

    \multirow{9}{=}{%
      \begin{minipage}[t]{7cm}
        \textbf{Bloque 3: Memoria y presentación oral}

        \vspace{1mm}
        RATFM3. Transmitir la información y los resultados de un proyecto de ingeniería de telecomunicación de manera oral y escrita.

        \vspace{1mm}
        RATFM8. Elaborar informes y memorias de calidad científico-tecnológica, que describan de forma clara y estructurada, un proyecto de ingeniería de telecomunicación, las referencias bibliográficas necesarias, una valoración de los resultados y una propuesta de mejoras.

        \vspace{1mm}
        RATFM9. Presentar y defender un proyecto en el ámbito de la Ingeniería de telecomunicación.
      \end{minipage}}
      & CE10. Expone oralmente los contenidos del trabajo con precisión y claridad y en el orden adecuado para facilitar su comprensión & & & \multirow[t]{9}{=}{} \\ \cline{2-4}
      & CE11. Utiliza en la presentación oral el vocabulario adecuado en cada circunstancia haciendo uso adecuado del léxico técnico cuando es necesario. & & & \\ \cline{2-4}
      & CE12. Muestra empatía con la audiencia transmitiendo seguridad en lo que se dice utilizando el tono y volumen adecuados. & & & \\ \cline{2-4}
      & CE13. Realiza la exposición del trabajo en el tiempo disponible para ello. & & & \\ \cline{2-4}
      & CE14. La memoria sigue la estructura esperada en un proyecto de ingeniería teniendo en cuenta las directrices marcadas por la normativa de TFM vigente. & & & \\ \cline{2-4}
      & CE15. La memoria está escrita tratando los diferentes temas con la profundidad necesaria y referenciando adecuadamente las fuentes de información. & & & \\ \cline{2-4}
      & CE16. Cuando la memoria describe un sistema, producto, prototipo, ... realizado, el contenido de ésta coincide con el sistema, producto, prototipo,… desarrollado. & & & \\ \cline{2-4}
      & CE17. La memoria está escrita con claridad, utilizando vocabulario adecuado y con un buen uso de las reglas ortográficas y gramaticales. & & & \\ \cline{2-4}
      & \cellcolor[HTML]{C0C0C0}\textbf{Calificación máxima 3 puntos} & & & \\ \cline{1-5}
  \end{tabularx}
\end{table}

\vspace{0.4cm}

\begin{center}
  \small
  \begin{tabularx}{\linewidth}{|X|}
    \hline
    \centering\textbf{Informe del Tribunal}\tabularnewline
    \hline
    \rule{0pt}{3.3cm} \\
    \hline
  \end{tabularx}
\end{center}

\end{landscape}
